\documentclass[border=4pt]{standalone}
\usepackage{amsmath,amssymb,mathtools,xcolor,varwidth}
\begin{document}
\begin{varwidth}{28em}
\large\bfseries
Now assemble the master key. The versed sine gives $F \propto QR/t^2$, while the swept area sets      $t \propto SP\cdot QT$; substitute and the force becomes reciprocally as the force-measure      solid $SP^2\cdot QT^2/QR$, taken in the limit as $Q\to P$ (Cor. 4 of Lem. X), and this solid is the heart      of the whole method. For ANY curve, pick $P$ and a nearby $Q$, draw $QR\parallel SP$ and $QT\perp SP$, and      compute this limit — it is the single most used formula in all of Book I. Props. 7, 11, and 13 will feed      the circle, the ellipse, and the parabola into exactly this construction. \textit{Q.E.D.}
\end{varwidth}
\end{document}
